% ============================================================
% MCM/ICM Summary Sheet Template (English)
% ============================================================
% COMAP Requirement:
%   The Summary Sheet is a single-page summary that precedes the
%   main paper. It does NOT count toward the 25-page limit.
%
% Usage:
%   In main.tex, after \begin{document}, add % ============================================================
% MCM/ICM Summary Sheet Template (English)
% ============================================================
% COMAP Requirement:
%   The Summary Sheet is a single-page summary that precedes the
%   main paper. It does NOT count toward the 25-page limit.
%
% Usage:
%   In main.tex, after \begin{document}, add % ============================================================
% MCM/ICM Summary Sheet Template (English)
% ============================================================
% COMAP Requirement:
%   The Summary Sheet is a single-page summary that precedes the
%   main paper. It does NOT count toward the 25-page limit.
%
% Usage:
%   In main.tex, after \begin{document}, add % ============================================================
% MCM/ICM Summary Sheet Template (English)
% ============================================================
% COMAP Requirement:
%   The Summary Sheet is a single-page summary that precedes the
%   main paper. It does NOT count toward the 25-page limit.
%
% Usage:
%   In main.tex, after \begin{document}, add \input{summary_sheet}
%   then \newpage before the main body begins.
%
% Six-Section Structure (R9 optimization, supersedes R1-R6 Chinese version):
%   1. Problem       -- concise problem restatement
%   2. Assumptions   -- key assumptions with brief justification
%   3. Model         -- model overview (no derivation details)
%   4. Results       -- headline numerical results
%   5. Conclusions   -- actionable conclusions
%   6. Sensitivity   -- sensitivity-analysis finding
%
% Single page limit: keep total length under ~500 words to fit 1 page.
% ============================================================

\thispagestyle{empty}
\begin{center}
\vspace*{0.5cm}
{\Large\textbf{Summary Sheet}} \\[0.3cm]
{\large Team \# [TEAM CONTROL NUMBER] \quad Problem [A/B/C/D/E/F]}
\vspace{0.5cm}
\end{center}

\rule{\textwidth}{0.5pt}

\begin{description}[leftmargin=0pt, itemindent=0pt, labelsep=0.5em, font=\bfseries]

% ============================================================
% Section 1: Problem
% ============================================================
\item[Problem.]
Restate the client's problem in one or two short sentences using
plain English (no jargon, no equations). State what decision or
insight the client needs and the time horizon.

\textit{Template:} The [client] needs to [decision/insight] in order
to [goal] within [time horizon].

% ============================================================
% Section 2: Assumptions
% ============================================================
\item[Assumptions.]
List the two or three most consequential assumptions. Each must be
followed by a one-clause justification (why it is reasonable, or
what evidence supports it).

\begin{enumerate}[leftmargin=1.5em, itemsep=0pt]
    \item \textbf{[Assumption 1].} \textit{Justification:} [brief
    reason or supporting reference].
    \item \textbf{[Assumption 2].} \textit{Justification:} [brief
    reason].
    \item \textbf{[Assumption 3].} \textit{Justification:} [brief
    reason].
\end{enumerate}

% ============================================================
% Section 3: Model
% ============================================================
\item[Model.]
Name the primary modeling approach (e.g., ``a system-dynamics model
coupled with Monte Carlo simulation'') and state the key state
variables or decision variables in one sentence. Do \emph{not} include
full equations here -- those belong in the main paper.

\textit{Template:} We construct a [model type] with [state/decision
variables] over [spatial/temporal scope], solved via [method].

% ============================================================
% Section 4: Results
% ============================================================
\item[Results.]
Report the headline numerical findings. Use specific numbers with
units and uncertainty ranges where applicable. Avoid vague phrases
like ``good performance'' -- the judges will look for concrete results.

\begin{itemize}[leftmargin=1.5em, itemsep=0pt]
    \item \textbf{[Result 1]:} [number $\pm$ uncertainty] [unit],
    [comparison to baseline if any].
    \item \textbf{[Result 2]:} [number] [unit], representing [%
    change] over [reference].
    \item \textbf{[Result 3]:} [number] [unit], validated against
    [validation method].
\end{itemize}

% ============================================================
% Section 5: Conclusions
% ============================================================
\item[Conclusions.]
State the actionable conclusion the client should draw. If the
problem is an F-problem, this is the headline recommendation that
the Memo to Client will elaborate. Be specific: ``do X to achieve
Y'' beats ``X is recommended''.

\textit{Template:} We recommend that [client] [specific action] to
[achieve specific outcome] within [timeframe].

% ============================================================
% Section 6: Sensitivity
% ============================================================
\item[Sensitivity.]
Identify the most sensitive parameter(s) and the direction of
sensitivity. State whether the conclusion is robust (holds across
plausible parameter ranges) or fragile (reverses for some values).

\textit{Template:} The conclusion is [robust / fragile] to
[parameter X] in the range [lo, hi]; [parameter Y] is the most
influential, with [direction] sensitivity.

\end{description}

\vspace{0.3cm}
\rule{\textwidth}{0.5pt}

\begin{center}
\textit{Keywords:} [keyword 1]; [keyword 2]; [keyword 3]; [keyword 4]
\end{center}

\vfill
\begin{flushright}
\textit{Summary Sheet does NOT count toward the 25-page limit.}
\end{flushright}

%   then \newpage before the main body begins.
%
% Six-Section Structure (R9 optimization, supersedes R1-R6 Chinese version):
%   1. Problem       -- concise problem restatement
%   2. Assumptions   -- key assumptions with brief justification
%   3. Model         -- model overview (no derivation details)
%   4. Results       -- headline numerical results
%   5. Conclusions   -- actionable conclusions
%   6. Sensitivity   -- sensitivity-analysis finding
%
% Single page limit: keep total length under ~500 words to fit 1 page.
% ============================================================

\thispagestyle{empty}
\begin{center}
\vspace*{0.5cm}
{\Large\textbf{Summary Sheet}} \\[0.3cm]
{\large Team \# [TEAM CONTROL NUMBER] \quad Problem [A/B/C/D/E/F]}
\vspace{0.5cm}
\end{center}

\rule{\textwidth}{0.5pt}

\begin{description}[leftmargin=0pt, itemindent=0pt, labelsep=0.5em, font=\bfseries]

% ============================================================
% Section 1: Problem
% ============================================================
\item[Problem.]
Restate the client's problem in one or two short sentences using
plain English (no jargon, no equations). State what decision or
insight the client needs and the time horizon.

\textit{Template:} The [client] needs to [decision/insight] in order
to [goal] within [time horizon].

% ============================================================
% Section 2: Assumptions
% ============================================================
\item[Assumptions.]
List the two or three most consequential assumptions. Each must be
followed by a one-clause justification (why it is reasonable, or
what evidence supports it).

\begin{enumerate}[leftmargin=1.5em, itemsep=0pt]
    \item \textbf{[Assumption 1].} \textit{Justification:} [brief
    reason or supporting reference].
    \item \textbf{[Assumption 2].} \textit{Justification:} [brief
    reason].
    \item \textbf{[Assumption 3].} \textit{Justification:} [brief
    reason].
\end{enumerate}

% ============================================================
% Section 3: Model
% ============================================================
\item[Model.]
Name the primary modeling approach (e.g., ``a system-dynamics model
coupled with Monte Carlo simulation'') and state the key state
variables or decision variables in one sentence. Do \emph{not} include
full equations here -- those belong in the main paper.

\textit{Template:} We construct a [model type] with [state/decision
variables] over [spatial/temporal scope], solved via [method].

% ============================================================
% Section 4: Results
% ============================================================
\item[Results.]
Report the headline numerical findings. Use specific numbers with
units and uncertainty ranges where applicable. Avoid vague phrases
like ``good performance'' -- the judges will look for concrete results.

\begin{itemize}[leftmargin=1.5em, itemsep=0pt]
    \item \textbf{[Result 1]:} [number $\pm$ uncertainty] [unit],
    [comparison to baseline if any].
    \item \textbf{[Result 2]:} [number] [unit], representing [%
    change] over [reference].
    \item \textbf{[Result 3]:} [number] [unit], validated against
    [validation method].
\end{itemize}

% ============================================================
% Section 5: Conclusions
% ============================================================
\item[Conclusions.]
State the actionable conclusion the client should draw. If the
problem is an F-problem, this is the headline recommendation that
the Memo to Client will elaborate. Be specific: ``do X to achieve
Y'' beats ``X is recommended''.

\textit{Template:} We recommend that [client] [specific action] to
[achieve specific outcome] within [timeframe].

% ============================================================
% Section 6: Sensitivity
% ============================================================
\item[Sensitivity.]
Identify the most sensitive parameter(s) and the direction of
sensitivity. State whether the conclusion is robust (holds across
plausible parameter ranges) or fragile (reverses for some values).

\textit{Template:} The conclusion is [robust / fragile] to
[parameter X] in the range [lo, hi]; [parameter Y] is the most
influential, with [direction] sensitivity.

\end{description}

\vspace{0.3cm}
\rule{\textwidth}{0.5pt}

\begin{center}
\textit{Keywords:} [keyword 1]; [keyword 2]; [keyword 3]; [keyword 4]
\end{center}

\vfill
\begin{flushright}
\textit{Summary Sheet does NOT count toward the 25-page limit.}
\end{flushright}

%   then \newpage before the main body begins.
%
% Six-Section Structure (R9 optimization, supersedes R1-R6 Chinese version):
%   1. Problem       -- concise problem restatement
%   2. Assumptions   -- key assumptions with brief justification
%   3. Model         -- model overview (no derivation details)
%   4. Results       -- headline numerical results
%   5. Conclusions   -- actionable conclusions
%   6. Sensitivity   -- sensitivity-analysis finding
%
% Single page limit: keep total length under ~500 words to fit 1 page.
% ============================================================

\thispagestyle{empty}
\begin{center}
\vspace*{0.5cm}
{\Large\textbf{Summary Sheet}} \\[0.3cm]
{\large Team \# [TEAM CONTROL NUMBER] \quad Problem [A/B/C/D/E/F]}
\vspace{0.5cm}
\end{center}

\rule{\textwidth}{0.5pt}

\begin{description}[leftmargin=0pt, itemindent=0pt, labelsep=0.5em, font=\bfseries]

% ============================================================
% Section 1: Problem
% ============================================================
\item[Problem.]
Restate the client's problem in one or two short sentences using
plain English (no jargon, no equations). State what decision or
insight the client needs and the time horizon.

\textit{Template:} The [client] needs to [decision/insight] in order
to [goal] within [time horizon].

% ============================================================
% Section 2: Assumptions
% ============================================================
\item[Assumptions.]
List the two or three most consequential assumptions. Each must be
followed by a one-clause justification (why it is reasonable, or
what evidence supports it).

\begin{enumerate}[leftmargin=1.5em, itemsep=0pt]
    \item \textbf{[Assumption 1].} \textit{Justification:} [brief
    reason or supporting reference].
    \item \textbf{[Assumption 2].} \textit{Justification:} [brief
    reason].
    \item \textbf{[Assumption 3].} \textit{Justification:} [brief
    reason].
\end{enumerate}

% ============================================================
% Section 3: Model
% ============================================================
\item[Model.]
Name the primary modeling approach (e.g., ``a system-dynamics model
coupled with Monte Carlo simulation'') and state the key state
variables or decision variables in one sentence. Do \emph{not} include
full equations here -- those belong in the main paper.

\textit{Template:} We construct a [model type] with [state/decision
variables] over [spatial/temporal scope], solved via [method].

% ============================================================
% Section 4: Results
% ============================================================
\item[Results.]
Report the headline numerical findings. Use specific numbers with
units and uncertainty ranges where applicable. Avoid vague phrases
like ``good performance'' -- the judges will look for concrete results.

\begin{itemize}[leftmargin=1.5em, itemsep=0pt]
    \item \textbf{[Result 1]:} [number $\pm$ uncertainty] [unit],
    [comparison to baseline if any].
    \item \textbf{[Result 2]:} [number] [unit], representing [%
    change] over [reference].
    \item \textbf{[Result 3]:} [number] [unit], validated against
    [validation method].
\end{itemize}

% ============================================================
% Section 5: Conclusions
% ============================================================
\item[Conclusions.]
State the actionable conclusion the client should draw. If the
problem is an F-problem, this is the headline recommendation that
the Memo to Client will elaborate. Be specific: ``do X to achieve
Y'' beats ``X is recommended''.

\textit{Template:} We recommend that [client] [specific action] to
[achieve specific outcome] within [timeframe].

% ============================================================
% Section 6: Sensitivity
% ============================================================
\item[Sensitivity.]
Identify the most sensitive parameter(s) and the direction of
sensitivity. State whether the conclusion is robust (holds across
plausible parameter ranges) or fragile (reverses for some values).

\textit{Template:} The conclusion is [robust / fragile] to
[parameter X] in the range [lo, hi]; [parameter Y] is the most
influential, with [direction] sensitivity.

\end{description}

\vspace{0.3cm}
\rule{\textwidth}{0.5pt}

\begin{center}
\textit{Keywords:} [keyword 1]; [keyword 2]; [keyword 3]; [keyword 4]
\end{center}

\vfill
\begin{flushright}
\textit{Summary Sheet does NOT count toward the 25-page limit.}
\end{flushright}

%   then \newpage before the main body begins.
%
% Six-Section Structure (R9 optimization, supersedes R1-R6 Chinese version):
%   1. Problem       -- concise problem restatement
%   2. Assumptions   -- key assumptions with brief justification
%   3. Model         -- model overview (no derivation details)
%   4. Results       -- headline numerical results
%   5. Conclusions   -- actionable conclusions
%   6. Sensitivity   -- sensitivity-analysis finding
%
% Single page limit: keep total length under ~500 words to fit 1 page.
% ============================================================

\thispagestyle{empty}
\begin{center}
\vspace*{0.5cm}
{\Large\textbf{Summary Sheet}} \\[0.3cm]
{\large Team \# [TEAM CONTROL NUMBER] \quad Problem [A/B/C/D/E/F]}
\vspace{0.5cm}
\end{center}

\rule{\textwidth}{0.5pt}

\begin{description}[leftmargin=0pt, itemindent=0pt, labelsep=0.5em, font=\bfseries]

% ============================================================
% Section 1: Problem
% ============================================================
\item[Problem.]
Restate the client's problem in one or two short sentences using
plain English (no jargon, no equations). State what decision or
insight the client needs and the time horizon.

\textit{Template:} The [client] needs to [decision/insight] in order
to [goal] within [time horizon].

% ============================================================
% Section 2: Assumptions
% ============================================================
\item[Assumptions.]
List the two or three most consequential assumptions. Each must be
followed by a one-clause justification (why it is reasonable, or
what evidence supports it).

\begin{enumerate}[leftmargin=1.5em, itemsep=0pt]
    \item \textbf{[Assumption 1].} \textit{Justification:} [brief
    reason or supporting reference].
    \item \textbf{[Assumption 2].} \textit{Justification:} [brief
    reason].
    \item \textbf{[Assumption 3].} \textit{Justification:} [brief
    reason].
\end{enumerate}

% ============================================================
% Section 3: Model
% ============================================================
\item[Model.]
Name the primary modeling approach (e.g., ``a system-dynamics model
coupled with Monte Carlo simulation'') and state the key state
variables or decision variables in one sentence. Do \emph{not} include
full equations here -- those belong in the main paper.

\textit{Template:} We construct a [model type] with [state/decision
variables] over [spatial/temporal scope], solved via [method].

% ============================================================
% Section 4: Results
% ============================================================
\item[Results.]
Report the headline numerical findings. Use specific numbers with
units and uncertainty ranges where applicable. Avoid vague phrases
like ``good performance'' -- the judges will look for concrete results.

\begin{itemize}[leftmargin=1.5em, itemsep=0pt]
    \item \textbf{[Result 1]:} [number $\pm$ uncertainty] [unit],
    [comparison to baseline if any].
    \item \textbf{[Result 2]:} [number] [unit], representing [%
    change] over [reference].
    \item \textbf{[Result 3]:} [number] [unit], validated against
    [validation method].
\end{itemize}

% ============================================================
% Section 5: Conclusions
% ============================================================
\item[Conclusions.]
State the actionable conclusion the client should draw. If the
problem is an F-problem, this is the headline recommendation that
the Memo to Client will elaborate. Be specific: ``do X to achieve
Y'' beats ``X is recommended''.

\textit{Template:} We recommend that [client] [specific action] to
[achieve specific outcome] within [timeframe].

% ============================================================
% Section 6: Sensitivity
% ============================================================
\item[Sensitivity.]
Identify the most sensitive parameter(s) and the direction of
sensitivity. State whether the conclusion is robust (holds across
plausible parameter ranges) or fragile (reverses for some values).

\textit{Template:} The conclusion is [robust / fragile] to
[parameter X] in the range [lo, hi]; [parameter Y] is the most
influential, with [direction] sensitivity.

\end{description}

\vspace{0.3cm}
\rule{\textwidth}{0.5pt}

\begin{center}
\textit{Keywords:} [keyword 1]; [keyword 2]; [keyword 3]; [keyword 4]
\end{center}

\vfill
\begin{flushright}
\textit{Summary Sheet does NOT count toward the 25-page limit.}
\end{flushright}
